% ============================================================
% 问题一的建模与求解
% Writer 为每个子问题创建独立文件（03-q1.tex, 04-q2.tex, ...）
% 每个文件包含完整的四段结构：分析→建立→求解→结果
% ============================================================
\section{问题一的建模与求解}

\subsection{问题分析}

（针对问题一，分析核心矛盾和解题切入点。
使用"我们"作为主语，如"针对问题一，我们注意到..."。
聚焦于：这个问题的主要矛盾是什么？我们打算用什么思路解决？
控制在 1--2 段。）

\subsection{模型建立}

（先说建模动机：为什么选择这个方法而不是其他方法？
观察到了什么现象或数据特征？
然后给出数学公式推导。公式之间用文字衔接，解释每一步的物理/数学含义。
避免连续堆砌公式。）

\subsection{模型求解}

（说明求解算法和计算过程。
图表穿插在首次引用处，使用 \verb|\begin{figure}[H]|。
每张图/表后至少 3 句分析文字。）

% 图表穿插示例（Writer 替换为实际图表）：
% \begin{figure}[H]
%   \centering
%   \includegraphics[width=0.8\textwidth]{figure/fig_q1_example.png}
%   \caption{问题一某某结果示意图}
%   \label{fig:q1_example}
% \end{figure}
%
% 如图~\ref{fig:q1_example} 所示，可以观察到...
% 这一现象表明...
% 特别值得注意的是...

\subsection{结果分析}

（结果用具体数值支撑每一个结论。
如："模型的 $R^2$ 达到 0.943，AIC 为 2408，较基准模型降低了 12.8\%。"
包括误差分析、灵敏度分析（如适用）。
结果对比优先使用表格呈现。）

% 结果表格示例（Writer 替换为实际数据）：
% \begin{table}[H]
% \centering
% \caption{问题一模型对比结果}
% \begin{tabular}{lccc}
% \toprule
% 模型 & $R^2$ & RMSE & AIC \\
% \midrule
% 基准模型 & 0.813 & 15.2 & 2762 \\
% 改进模型 & 0.943 & 8.7 & 2408 \\
% \bottomrule
% \end{tabular}
% \end{table}
